\begin{examplebox}
\textbf{\textcolor{CExample}{例 1（基础）}}

\textbf{\textcolor{CExample}{题目：}} 在等差数列 $\{a_n\}$ 中，已知 $a_2 + a_5 = 10$，求 $a_3 + a_4$ 的值。

\textbf{\textcolor{CExample}{解题步骤：}}
\begin{description}[style=nextline, leftmargin=2.8em, labelsep=0.5em, itemsep=0.45em]
    \item[\textbf{第一步（找下标和）：}] 观察下标：$2+5 = 7$，$3+4 = 7$，满足 $m+n = p+q$。
    \par\textbf{理由：} 验证是否满足结论条件。
    \item[\textbf{第二步（套用结论）：}] 由结论 $a_m+a_n = a_p+a_q$，可得
    \[
    a_2 + a_5 = a_3 + a_4.
    \]
    \par\textbf{理由：} 下标和相等，项和相等。
    \item[\textbf{第三步（计算结果）：}] 因为 $a_2+a_5 = 10$，所以 $a_3 + a_4 = 10$。
    \par\textbf{理由：} 直接代入。
\end{description}

\textbf{\textcolor{CExample}{关键结论：}} \textbf{无需知道公差和首项，直接由下标和相等快速口算。}

\medskip
\textbf{\textcolor{CExample}{例 2（稍变形）}}

\textbf{\textcolor{CExample}{题目：}} 在等差数列 $\{a_n\}$ 中，已知 $a_1 + a_4 + a_7 = 15$，求 $a_4$ 的值。

\textbf{\textcolor{CExample}{解题步骤：}}
\begin{description}[style=nextline, leftmargin=2.8em, labelsep=0.5em, itemsep=0.45em]
    \item[\textbf{第一步（构造下标和等）：}] 注意到 $1+7 = 8$，而 $2\times 4 = 8$，即 $1+7 = 4+4$。
    \par\textbf{理由：} 将三项和中的 $a_1,a_7$ 组合，利用中项推广。
    \item[\textbf{第二步（应用中项推广）：}] 由推广结论，$a_1 + a_7 = 2a_4$。
    \par\textbf{理由：} 当下标满足 $m+n=2k$ 时，两项和等于中项的两倍。
    \item[\textbf{第三步（代入求和）：}] 原式变为
    \[
    a_1 + a_4 + a_7 = (a_1 + a_7) + a_4 = 2a_4 + a_4 = 3a_4.
    \]
    \par\textbf{理由：} 结合已知求和。
    \item[\textbf{第四步（解方程）：}] 由 $3a_4 = 15$，解得 $a_4 = 5$。
    \par\textbf{理由：} 直接解方程。
\end{description}

\textbf{\textcolor{CExample}{关键结论：}} \textbf{下标对称时可先转化为中项关系，再整体求解。}
\end{examplebox}
